% from ega4-13-fr.tex
\oldpage[IV-3]{190}
$B=A[T_1,\ldots,T_n]/\mathfrak J$. Soit $(\mathfrak J_\lambda)$ la famille des
idéaux de type fini de $A[T_1,\ldots,T_n]$ contenus dans $\mathfrak J$, de
sorte que $\mathfrak J$ est la réunion filtrante des $\mathfrak J_\lambda$~;
si $X$ et les
$X_\lambda=\operatorname{Spec}(A[T_1,\ldots,T_n]/\mathfrak J_\lambda)$ sont
considérés comme des sous-préschémas fermés de
$Z=\operatorname{Spec}(A[T_1,\ldots,T_n])$, on a donc, pour les espaces
sous-jacents, $X=\bigcap_\lambda X_\lambda$. Si $f_\lambda:X_\lambda\to S$ est
le morphisme structural, on en déduit que
$f^{-1}(s)=\bigcap_\lambda f_\lambda^{-1}(s)$ pour tout $s\in S$, et comme les
ensembles $f_\lambda^{-1}(s)$ sont fermés dans l'espace noethérien $Z_s$, il
existe un $\lambda$ (dépendant de $s$) tel que
$f^{-1}(s)=f_\lambda^{-1}(s)$. Si alors, pour tout $x\in X$, on pose
$d(x)=\dim_x(f^{-1}(f(x)))$, $d_\lambda(x)=\dim_x(f_\lambda^{-1}(f_\lambda(x)))$,
ce qui précède prouve que l'on a $d(x)=\inf_\lambda d_\lambda(x)$~; les
fonctions $d_\lambda$ étant semi-continues supérieurement d'après la première
partie de la démonstration, il en est de même de $d$. C.Q.F.D.
\end{proof}

\begin{corollary}[13.1.4]
\EGAReaderTarget{ega4.c04.cor.13-1-4}{IV.13.1.4-fr}
Sous les hypothèses de (13.1.3), l'ensemble des $x\in X$ tels que $x$ soit
isolé dans $f^{-1}(f(x))$ est ouvert dans $X$.
\end{corollary}

\begin{proof}
En effet, c'est le complémentaire de $F_1(X)$ ($\mathbf 0$, 14.1.10).
\end{proof}

On notera qu'on retrouve ainsi, sous des hypothèses plus générales, la
conséquence (\textbf{III}, 4.4.10) du «~Main theorem~» de Zariski.

\begin{corollary}[13.1.5]
\label{IV.13.1.5-fr}
Sous les hypothèses de (13.1.3), supposons en outre que $f$ soit un morphisme
fermé. Alors, pour tout entier $n$, l'ensemble des $y\in Y$ tels que
$\dim(f^{-1}(y))\geq n$ est fermé, autrement dit, l'application
$y\mapsto\dim(f^{-1}(y))$ est semi-continue supérieurement~; en particulier,
si $y$ est spécialisation de $y'$, on a
$\dim(f^{-1}(y))\geq\dim(f^{-1}(y'))$.
\end{corollary}

\begin{proof}
En effet, dire que $\dim(f^{-1}(y))\geq n$ signifie que $y\in f(F_n(X))$
($\mathbf 0$, 14.1.6).
\end{proof}

\begin{corollary}[13.1.6]
\label{IV.13.1.6-fr}
Soient $X$, $Y$ deux préschémas irréductibles, $f:X\to Y$ un morphisme
dominant localement de type fini. Soit $\xi$ (resp. $\eta$) le point générique
de $X$ (resp. $Y$) et soit $e=\dim(f^{-1}(\eta))$. Alors, pour tout $x\in X$,
on a $\dim_x(f^{-1}(f(x)))\geq e$.
\end{corollary}

\begin{proof}
En effet, l'ensemble des $x$ tels que $\dim_x(f^{-1}(f(x)))<e$ est ouvert en
vertu de (13.1.3), et comme il ne peut contenir $\xi$, il est vide.
\end{proof}

Notons enfin le résultat plus facile suivant~:

\begin{proposition}[13.1.7]
\label{IV.13.1.7-fr}
Soient $Y$ un préschéma quasi-compact, $f:X\to Y$ un morphisme de type fini.
Il existe un entier $n$ tel que, pour tout $y\in Y$, on ait
$\dim(f^{-1}(y))\leq n$.
\end{proposition}

\begin{proof}
Comme il y a un recouvrement ouvert affine fini $(V_i)$ de $Y$ tel que chaque
$f^{-1}(V_i)$ soit réunion finie d'ouverts affines, on est aussitôt ramené au
cas où $Y=\operatorname{Spec}(A)$ et $X=\operatorname{Spec}(B)$ sont affines,
$B$ étant une $A$-algèbre de type fini. Si $B$ admet un système de $n$
générateurs, alors, pour tout $y\in Y$, $B\otimes_Ak(y)$ est une
$k(y)$-algèbre admettant $n$ générateurs, donc $\dim(f^{-1}(y))\leq n$ en
vertu de (4.1.1).
\end{proof}

\subsection{Morphismes équidimensionnels : cas des morphismes dominants de préschémas irréductibles}
\label{subsection:IV.13.2-fr}

\begin{env}[13.2.1]
\label{IV.13.2.1-fr}
Soient $Y$ un préschéma irréductible, $X$ un préschéma irréductible,
$f:X\to Y$ un morphisme dominant localement de type fini~; soit $\eta$ le
point générique de $Y$. On sait (13.1.6) que pour tout $x\in X$, on a
$\dim_x(f^{-1}(f(x)))\geq\dim(f^{-1}(\eta))$.
\end{env}

